\documentclass[sigconf,anonymous]{acmart}

\usepackage{amsthm}
\usepackage{booktabs}
\usepackage{graphicx}
\usepackage{placeins}
\usepackage{balance}

\newtheorem{proposition}{Proposition}
\newtheorem{corollary}{Corollary}

% Remove ACM-specific metadata for preprint
\settopmatter{printacmref=false}
\renewcommand\footnotetextcopyrightpermission[1]{}
\pagestyle{plain}

\title{A Structural Threshold in Decision Capacity Governs Collapse in Self-Play Reinforcement Learning}

\author{Anonymous}
\affiliation{\institution{Anonymous}}
\email{anonymous@example.com}

\begin{abstract}
We show that a threshold in decision capacity determines whether self-play reinforcement learning agents collapse under asymmetric rule perturbations. Across poker variants (Kuhn, Leduc, Leduc-4), matrix games (Matching Pennies), a dice game (Liar's Dice; 24{,}576 info sets), and four learning algorithms (Q-Learning, SARSA, REINFORCE, DQN), eliminating all \emph{strategically meaningful} decisions causes rapid convergence to a deterministic exploitation attractor (DEA)---a fixed point at near-maximal loss. Preserving even a single contingent decision point prevents this collapse. A frozen baseline and fixed-opponent control confirm the mechanism is co-adaptation under constraint, not the perturbation itself. The phenomenon is timing-invariant, fully reversible upon action restoration, and intensifies under function approximation. These results establish a sharp, structural threshold between zero and minimal decision capacity that governs exploitability in self-play multi-agent systems.
\end{abstract}

\keywords{Multi-agent reinforcement learning; self-play; robustness; game theory; action-space perturbation}

\begin{document}
\maketitle

% ===================================================================
\section{Introduction}

Multi-agent reinforcement learning (MARL) agents trained through self-play have achieved superhuman performance in complex games~\cite{silver2018general,brown2019superhuman,vinyals2019grandmaster}, yet their robustness to structural changes in the environment remains poorly understood. Prior work has focused primarily on adversarial perturbations to observations or rewards~\cite{gleave2020adversarial}, opponent modelling under distribution shift~\cite{foerster2018learning}, or training stability in population-based methods~\cite{balduzzi2019open}. Structural changes to the \emph{action space}---where an agent permanently loses access to certain actions---remain largely unexplored.

Such perturbations arise naturally in practice. In robotics, hardware failures may disable actuators, eliminating actions mid-deployment. In financial trading, regulatory changes can restrict previously available strategies. In multi-agent software systems, API deprecations may remove action endpoints. Understanding how self-play agents respond to these asymmetric capability losses is essential for deploying RL systems where the action space is not guaranteed to remain static.

We study this question in imperfect-information games by removing one player's ability to bet or raise at subsets of decision nodes. This creates a controlled perturbation that reduces the agent's \emph{contingent action capacity} (CAC)---the number of information sets at which it retains more than one legal action.

Our central finding is a sharp threshold effect. When CAC drops to zero (every decision is forced), adaptive self-play agents converge to a \emph{deterministic exploitation attractor} (DEA)---a fixed point characterised by near-maximal loss and near-zero reward variance. When even a single decision point is preserved (CAC $\geq 1$), agents stabilise near Nash equilibrium. A frozen baseline and fixed-opponent control isolate continued co-adaptation as the mechanism. Our contributions:
\begin{itemize}
  \item We identify a sharp CAC threshold governing collapse in self-play RL and provide formal propositions characterising the zero-contingency fixed point.
  \item We isolate co-adaptation under constraint as the mechanism via frozen baseline and fixed-opponent comparisons.
  \item We demonstrate that collapse persists and intensifies under function approximation (DQN).
  \item We replicate across five games (1--24{,}576 info sets), four algorithms, multiple perturbation schedules, and show full reversibility. Boundary conditions (IPD, Liar's Dice) sharpen the threshold definition.
\end{itemize}

% ===================================================================
\section{Related Work}

\paragraph{Solving imperfect-information games.}
CFR~\cite{zinkevich2007regret} and variants have solved increasingly large poker games~\cite{bowling2015heads,moravcik2017deepstack,brown2018superhuman,brown2019superhuman}. These systems assume a fixed game structure; we study what happens when it changes asymmetrically.

\paragraph{Self-play RL.}
From TD-Gammon~\cite{tesauro1994td} through AlphaZero~\cite{silver2018general} and AlphaStar~\cite{vinyals2019grandmaster}, self-play has driven major advances. NFSP~\cite{heinrich2016deep} and PSRO~\cite{lanctot2017unified} address stability, but self-play can cycle or exhibit non-transitivity~\cite{balduzzi2019open,lanctot2019openspiel}. We identify a distinct failure mode: co-adaptation-driven collapse under asymmetric action-space constraints.

\paragraph{MARL robustness.}
The literature encompasses adversarial policies~\cite{gleave2020adversarial}, opponent-learning awareness~\cite{foerster2018learning}, and distributional robustness~\cite{zhang2021multi}. We perturb the action space itself, revealing a qualitative threshold that continuous perturbations cannot produce.

\paragraph{Action masking and constrained RL.}
Action masking is standard in game AI~\cite{huang2022closer}, and constrained MDPs formalise restrictions~\cite{altman1999constrained}. The dynamic consequences of mid-training removal under self-play have not been studied.

\paragraph{Exploitability.}
The standard game-theoretic metric~\cite{johanson2007computing,timbers2022approximate}. We connect our threshold to the best-response structure of reduced games via formal propositions.

% ===================================================================
\section{Background}

\textbf{Kuhn Poker.} Three cards ($J\!<\!Q\!<\!K$), two actions, ante~1. Nash P0: $-1/18$. 12 info sets.
\textbf{Leduc Poker.} Six cards, three actions, two rounds. Nash P0: $\approx\!-0.087$. 288 info sets.
\textbf{Leduc-4.} Four ranks, three suits (12 cards). 504 info sets. Nash P0: $\approx\!-0.096$.
\textbf{Liar's Dice.} 1-die: 13 actions, 24{,}576 info sets. 2-dice: 25 actions, $200{,}000{+}$ info sets.
\textbf{Coordination Game.} Cooperative: 2 agents match a random target, 3 actions, 10 rounds.
\textbf{Negotiation.} Ultimatum: P0 offers a split (11 actions), P1 accepts/rejects, 5 rounds.
\textbf{Matching Pennies.} Single-shot, Nash: 50/50, value 0.

\textbf{Agents.} CFR (frozen Nash), Q-Learning (tabular $\varepsilon$-greedy), QL-Frozen (frozen at perturbation), DQN (2-layer MLP), SARSA, REINFORCE.

% ===================================================================
\section{Methodology}

Each experiment: 20 seeds, 20{,}000 episodes (50{,}000 for DQN). Perturbation at midpoint.

\textbf{CAC.} Number of info sets where the perturbed agent retains $>1$ legal action. Unweighted by reach probability (see Limitations).

\textbf{Self-play.} Single agent plays both roles; Q-values indexed by player-specific info states. Verified equivalent to separate agents (Appendix~\ref{app:separate}).

\textbf{Statistics.} Paired $t$-tests, bootstrap 95\% CIs (10{,}000 resamples), Cohen's $d$.

% ===================================================================
\section{Theory}

\begin{proposition}[Zero-contingency exploitation]
\label{prop:zero}
Let $G'$ be the reduced game with $\mathrm{CAC}(P_0)=0$. Then $v_0(G')$ equals the best-response value against $P_0$'s forced policy: $v_0(G') = -\max_{\pi_1} \sum_z u_1(z) \cdot \rho^{\sigma_0^f,\pi_1}(z)$.
\end{proposition}
\begin{proof}
$P_0$ has no choices; the game is a single-agent optimisation for $P_1$ with a unique pure BR. By the minimax theorem, $v_0(G')$ is determined uniquely.
\end{proof}

\begin{proposition}[Residual contingency bound]
\label{prop:residual}
If $|A(h^*)| \geq 2$ for some $h^* \in \mathcal{I}_0$, then $P_1$ cannot adopt the pure exploitation strategy of Prop.~\ref{prop:zero}, and $v_0(G') \geq v_0^{\mathrm{Nash}}(G) - \delta(h^*)$.
\end{proposition}
\begin{proof}
$P_0$'s choice at $h^*$ forces $P_1$ to mix or condition, preventing convergence to the zero-contingency point.
\end{proof}

\begin{corollary}
The $\mathrm{CAC}=1\to0$ transition qualitatively changes the BR structure: from strategic interaction to trivially computable pure exploitation.
\end{corollary}

% ===================================================================
\section{Experiments}

\subsection{The Phenomenon}

\begin{table}[t]
\centering
\caption{Kuhn, zero contingency (20 seeds).}
\label{tab:kuhn-full}
\begin{tabular}{lcccc}
\toprule
Agent & Pre & Post & $p$ & $d$ \\
\midrule
CFR & $-0.060$ & $-0.221$ & ${<}0.0001$ & $-6.9$ \\
Q-Learning & $-0.041$ & $\mathbf{-0.926}$ & ${<}0.0001$ & $-42.0$ \\
\bottomrule
\end{tabular}
\end{table}

Under zero contingency (Table~\ref{tab:kuhn-full}), Q-Learning converges to the DEA at $-0.926$ (normalized: $0.27$; $d=-42.0$) within four episodes. Under residual contingency (CAC$=1$), Q-Learning drops only to $-0.065$ ($d=-0.8$) and stabilises.

\subsection{The Threshold}

\begin{table}[t]
\centering
\caption{CAC sweep (Kuhn).}
\label{tab:capacity}
\begin{tabular}{clccc}
\toprule
CAC & & CFR & QL & Norm. \\
\midrule
0 & Zero-cont. & $-0.221$ & $\mathbf{-0.926}$ & $0.27$ \\
1 & Residual & $-0.053$ & $-0.065$ & $0.48$ \\
2 & Full & $-0.053$ & $-0.034$ & $0.49$ \\
\bottomrule
\end{tabular}
\end{table}

The $0{\to}1$ jump is $\Delta=0.86$; the $1{\to}2$ jump is $\Delta=0.03$---a sharp threshold, not a gradient (Table~\ref{tab:capacity}).

\subsection{The Mechanism}

\textbf{Frozen baseline.} QL-Frozen: $-0.141$ vs Q-Learning: $-0.927$ ($p=0.004$, $d=-2.6$). Co-adaptation, not the constraint, drives collapse.

\textbf{Fixed opponent.} Against static Nash: $-0.228$ (no collapse). Self-play: $-0.926$. Co-adaptation is necessary.

\subsection{Generalisation}

\begin{table}[t]
\centering
\caption{Algorithm invariance (Kuhn, 5 seeds).}
\label{tab:algorithms}
\begin{tabular}{lccc}
\toprule
Algorithm & Type & Post & $d$ \\
\midrule
Q-Learning & Tab. & $\mathbf{-0.927}$ & $-66.1$ \\
SARSA & Tab. & $\mathbf{-0.927}$ & $-66.1$ \\
REINFORCE & Tab. & $\mathbf{-0.500}$ & $-39.6$ \\
DQN & Neural & $\mathbf{-0.994}$ & $-24.4$ \\
\bottomrule
\end{tabular}
\end{table}

All collapse (Table~\ref{tab:algorithms}). DQN collapses hardest---policy entropy drops to near zero and Q-value gaps spike post-perturbation, confirming rapid deterministic convergence.

\begin{table}[t]
\centering
\caption{Cross-game collapse severity.}
\label{tab:cross-game}
\begin{tabular}{lccc}
\toprule
Game & QL Post & Info sets & Norm. \\
\midrule
Match.\ Pennies & $-0.851$ & 1 & $0.07$ \\
Kuhn & $-0.926$ & 12 & $0.27$ \\
Leduc & $-0.252$ & 288 & $0.49$ \\
Leduc-4 & $-0.185$ & 504 & $0.49$ \\
Liar's Dice (1d) & $-0.032^*$ & 24k & $0.48$ \\
Coordination & $+1.44^\ddagger$ & -- & -- \\
Liar's Dice (2d) & $+0.008^*$ & 200k+ & $0.50$ \\
\bottomrule
\end{tabular}

{\footnotesize $^*$Challenge-only retains strategic value. $^\ddagger$Cooperative; degradation without collapse.}
\end{table}

Collapse holds across poker and matrix games (Table~\ref{tab:cross-game}). Liar's Dice shows no collapse---a boundary condition.

\subsection{Boundary Conditions}
\label{sec:boundaries}

\textbf{IPD.} Removing ``cooperate'' pushes toward Nash (always-defect). No collapse.

\textbf{Liar's Dice.} Challenge-only retains strategic value (contingent on private die). No collapse at either scale (1d/2d). Confirms the threshold depends on \emph{strategic flexibility}, not action count.

\textbf{Non-zero-sum domains.} In the cooperative Coordination game, forcing P0 to a single action (CAC $= 0$) degrades performance ($p = 0.001$) but does not produce convergence to the DEA. In the Negotiation game, P1 shifts toward rejection but does not converge to unconditional rejection. Cooperative dynamics produce bounded degradation, not collapse---the threshold interacts with interaction structure.

\textbf{Timing.} Episodes 3k/10k/17k yield identical collapse ($-0.926$/$-0.927$/$-0.925$).

\subsection{Dynamics}

\textbf{Recovery.} Restoring actions produces full recovery ($\Delta=+0.90$) within four episodes---symmetric with collapse. The DEA is a maintained attractor, not corrupted representation.

\textbf{Exploitability.} Exact BR computation confirms post-perturbation exploitability spikes to maximum.

\textbf{Variance.} Post-collapse $V_{\mathrm{total}}=5{\times}10^{-6}$: the DEA is a deterministic fixed point.

% ===================================================================
\section{Conclusion}

We identified a sharp threshold in contingent action capacity governing collapse in self-play RL. The finding is game-general (1--200{,}000+ info sets), algorithm-invariant (including DQN), co-adaptation-driven, timing-invariant, and fully reversible. Boundary conditions (IPD, Liar's Dice) show the threshold depends on strategic flexibility, not action count. Entropy regularisation does not prevent collapse (Appendix~\ref{app:entropy}).

\textbf{Limitations.} CAC is unweighted by reach probability. Reversibility may not hold under deeper networks.

\textbf{Future work.} Larger games with neural CFR; cooperative settings; formal exploitability bounds; reach-weighted capacity.

% ===================================================================
\balance
\bibliographystyle{ACM-Reference-Format}
\bibliography{references}

% ===================================================================
\clearpage
\appendix

\section{Hyperparameter Sensitivity}
\label{app:hyperparam}

\begin{table}[h]
\centering
\begin{tabular}{lccc}
\toprule
& $\alpha\!=\!0.01$ & $\alpha\!=\!0.1$ & $\alpha\!=\!0.3$ \\
\midrule
$\varepsilon\!=\!0.05$ & $-0.974$ & $-0.976$ & $-0.976$ \\
$\varepsilon\!=\!0.15$ & $-0.924$ & $-0.927$ & $-0.927$ \\
$\varepsilon\!=\!0.30$ & $-0.849$ & $-0.853$ & $-0.851$ \\
\bottomrule
\end{tabular}
\caption{Collapse across all hyperparameter combinations.}
\end{table}

\section{Separate Self-Play Verification}
\label{app:separate}
Independent P0/P1 Q-tables: $-0.926\pm0.003$ vs shared: $-0.927$ (diff: $+0.0008$).

\section{Entropy Regularisation}
\label{app:entropy}

\begin{table}[h]
\centering
\begin{tabular}{lc}
\toprule
$\tau$ & QL Post \\
\midrule
$0.00$ & $-0.927$ \\
$0.05$ & $-0.925$ \\
$0.10$ & $-0.923$ \\
$0.20$ & $-0.926$ \\
\bottomrule
\end{tabular}
\caption{Entropy regularisation has no effect on collapse.}
\end{table}

\end{document}
